\documentclass[oneside]{scrbook} %KOMA-Script book
\usepackage[italian]{babel}
%Per i commenti
\usepackage{comment}
%Per l'interlinea
\usepackage{setspace}
%Codifica caratteri
\usepackage[utf8]{inputenc}
\usepackage{xspace}
\usepackage{eurosym}
%Fonts Scientifici
\usepackage{amsfonts}
\usepackage{amsmath}
\usepackage{amssymb}
\usepackage{textcomp}
%Pacchetto grafici
\usepackage{tikz}
\usepackage{pgfplots}
\pgfplotsset{/pgf/number format/use comma,compat=newest}
\usepackage{picture}
\usepackage{graphicx}
%Impostazioni tabelle avanzate e colorate
\usepackage{multicol}
\usepackage{multirow}
\usepackage{soul}
\usepackage{ulem}
\usepackage{booktabs}
%\usepackage{subfig}
\usepackage{subfigure}
\usepackage{array}
\usepackage{color}
\usepackage{colortbl}
%Pacchetto per gli pseudocodici
\usepackage{algpseudocode}
%\usepackage{listings}
\usepackage{listings,xcolor,courier,bookmark}
\usepackage{listingsutf8}
%Impostazioni tabelle avanzate
\usepackage{multicol}
\usepackage{multirow}
\usepackage{rotating}
\definecolor{darkblue}{named}{blue}
\definecolor{darkred}{named}{red}
\definecolor{grau}{named}{gray}
%Per i margini
\usepackage{chngpage}
\sloppy

\graphicspath{{Immagini/}}

\let\Righttorque\relax

\lstset{
	captionpos=b,
	commentstyle=\color[rgb]{0.133,0.545,0.133},
	keywordstyle=\color{darkblue},
	stringstyle=\color{purple},
	extendedchars=true,
	basicstyle=\small\ttfamily,
	showstringspaces=false,
	tabsize=2,
	numbers=left,
	numberstyle=\tiny,
	breakautoindent  = true,
	breakindent      = 2em,
	breaklines       = true,
	postbreak        = ,
	prebreak         = \raisebox{-.8ex}[0ex][0ex]{\Righttorque},
	showspaces=false,
	showtabs=false,
	showstringspaces=false,
	language=C++,
	frame=single,
	morecomment=[s]{°°}
}

\renewcommand*{\lstlistingname}{Codice}
\renewcommand*{\lstlistlistingname}{Elenco dei codici}

\usepackage{fancyhdr}
\pagestyle{fancy}
\renewcommand{\chaptermark}[1]{\markboth{\MakeUppercase{\thechapter.\ #1}}{}}
\fancypagestyle{plain} %stile delle pagine anche a quella iniziale del capitolo

\fancyhead{}
\fancyfoot{}

\fancyhead[R]{\bfseries \nouppercase{\leftmark}}
\fancyhead[L]{
	\includegraphics*[scale=0.8]{logo_dieti2.png}
}
\fancyfoot[R]{\thepage}
\fancyfoot[L]{Tesina di Impianti di Elaborazione}
\renewcommand{\headrulewidth}{0.4pt}
\renewcommand{\footrulewidth}{0.4pt}

\date{}
\cfoot{}
\usepackage{eso-pic,graphicx}

%Per sottosezioni
\setcounter{secnumdepth}{4}
\setcounter{tocdepth}{4}

\begin{document}

	% Per evitare che i collegamenti nell'indice abbiano i riquadri rossi
	\hypersetup {linkbordercolor=white}

	\AddToShipoutPicture{\AtPageCenter{\makebox(0,0){\includegraphics{logo.png}}}}
	\frontmatter
	\pagenumbering{Roman}

	\begin{titlepage}
		\centering
		%{\Huge \textsc{Università degli Studi di Napoli ``Federico II''}}

		%\vspace*{\stretch{2.5}}

		%\includegraphics[width=1\linewidth]{logo_federico_II.png}

		\includegraphics[width=1\linewidth]{logo_dieti2.png}

		\vspace*{2cm}

		{\Huge \textsl{Tesina di\\IMPIANTI DI ELABORAZIONE}\par}

		\vspace*{2cm}

		{\huge \text{Prof. Domenico Cotroneo}}

		\vspace*{5cm}

		{\LARGE \textit{Andrea Scognamiglio - Matr. M63/598\\Cristian Tommasino - Matr. M63/615}}

		\vspace*{2.5cm}

		{\Large \textsc{A.A. 2017/2018}}
	\end{titlepage}

	\setcounter{page}{1}

	\newpage
	\tableofcontents

	\newpage
	\mainmatter

	% \input{.tex}

\end{document}
